\section*{Chapitre \ref{ch:qm:high-frequency-econometrics} --- Économétrie haute fréquence}

\begin{solution}{exo:qm:high-frequency-econometrics:1}
$\E[\mathrm{RV}] = 10^{-4} + 2 \times 23\,400 \times (2 \times 10^{-4})^2 = 10^{-4} + 1.87 \times 10^{-3} = 1.97 \times 10^{-3}$~: une volatilité quotidienne de 4.4~\% pour une vraie volatilité de 1~\%.
\end{solution}

\begin{solution}{exo:qm:high-frequency-econometrics:2}
$2\sqrt{0.0004} = 0.04$~dollar, soit quatre cents.
\end{solution}

\begin{solution}{exo:qm:high-frequency-econometrics:3}
Le milieu ne saute pas entre le cours acheteur et le cours vendeur avec le signe de la transaction, donc le rebond disparaît~; mais il ne bouge que lorsque les cotations franchissent un point de la grille, si bien que c'est le prix efficient
arrondi aux demi-ticks, et l'erreur d'arrondi joue le rôle de bruit.
\end{solution}

\begin{solution}{exo:qm:high-frequency-econometrics:4}
$\mathrm{IQ} = (10^{-4})^2 = 10^{-8}$ et $\omega^4 = (4 \times 10^{-8})^2 = 1.6 \times 10^{-15}$, donc $n^* = (10^{-8}/6.4 \times 10^{-15})^{1/3} = 116$ rendements, soit un toutes les 202 secondes.
\end{solution}

\begin{solution}{exo:qm:high-frequency-econometrics:5}
Un rendement observé vaut $r_i + u_i - u_{i-1}$, de variance $\sigma^2\Delta + 2\omega^2$~; deux rendements consécutifs partagent $u_{i-1}$ avec des signes opposés, si bien que leur covariance vaut
$-\omega^2$. Pour l'action du chapitre à une seconde, $\sigma^2\Delta = 2.48 \times 10^{-4}/23\,400 = 1.06 \times 10^{-8}$ et $\omega^2 = 2.56 \times 10^{-8}$~: $-2.56/(1.06 + 5.12) =
-0.41$, comme mesuré dans la simulation ($-0.41$).
\end{solution}

\begin{solution}{exo:qm:high-frequency-econometrics:6}
Avec une négociation indépendante, $0.2 \times 1/15 = 1.3\%$ des intervalles d'une seconde contiennent une transaction dans les deux titres~; dans les autres, au moins un rendement au dernier tick est nul, ce qui explique pourquoi la
corrélation à une seconde s'effondre.
\end{solution}

\begin{solution}{exo:qm:high-frequency-econometrics:7}
\texttt{signature()} et \texttt{mse\_by\_interval()}~: transactions 60.1~\%, 35.1~\%, 26.1~\% et 25.4~\% à 1, 5, 60 et 300 secondes~; milieux 36.8~\% à une seconde~; prix efficient 25.0~\% partout. La
racine de l'erreur quadratique moyenne relative vaut 4.80 à une seconde et est minimale autour d'une minute (0.11 à 60 secondes, 0.12 à 120).
\end{solution}

\begin{solution}{exo:qm:high-frequency-econometrics:8}
Les rendements à une seconde de deux actions contiennent rarement des transactions dans les deux titres, si bien que la corrélation réalisée au dernier tick est biaisée vers zéro (l'\omterm{def:qm:high-frequency-econometrics:epps}{effet Epps})~: 0.12 est un artefact de
l'échantillonnage, et non une propriété de la paire. Utiliser Hayashi--Yoshida, ou échantillonner sur plusieurs minutes, ou un \omterm{def:qm:estimation:estimator}{estimateur} synchronisé robuste au bruit, puis comparer.
\end{solution}

\begin{solution}{pb:qm:high-frequency-econometrics:1}
\textbf{1.} 60.1~\%, 35.1~\%, 26.1~\% et 25.4~\%.
\textbf{2.} 36.8~\% à partir des milieux de cotation à une seconde~; 25.0~\% à partir du prix efficient pour chaque intervalle.
\textbf{3.} Estimé à partir de $\mathrm{RV}/2n$~: 1.75 points de base~; vrai~: 1.60 (une transaction se situe uniformément à moins d'un tick du prix efficient, un tick valant 2.8 points de base).
\textbf{4.} 1.14~cent, contre une fourchette réelle d'un cent.
\textbf{5.} Chaque rendement à une seconde porte une différence de bruit de variance $2\omega^2$, soit environ cinq fois la variance efficiente par seconde, et les 23\,400 rendements s'additionnent~: la RV vaut 5.8 fois
la \omterm{def:qm:volatility-models:rv}{variance intégrée}.
\textbf{6.} $\omega^2/\mathrm{IV} = 1.24 \times 10^{-4}$.
\textbf{7.} $n^* = 254$, soit un rendement toutes les 92 secondes (82 avec le vrai bruit).
\textbf{8.} 4.80 à une seconde, 0.98 à cinq, 0.18 à trente, 0.11 à soixante, 0.12 à 120, 0.14 à 300 et 0.34 à 1\,800 secondes.
\textbf{9.} Elle utilise 254 rendements sur 23\,400, soit environ 1~\%.
\textbf{10.} $n^* \propto \mathrm{IQ}^{1/3}$ et l'IQ varie avec le carré de la variance, si bien que doubler la variance multiplie $n^*$ par $2^{2/3} = 1.59$~: échantillonner environ toutes les 58 secondes.
\textbf{11.} 25.2~\%, 25.1~\% et 25.1~\%, avec des erreurs relatives quotidiennes de 11.9~\%, 7.2~\% et 4.8~\% (RV à cinq minutes~: 25.4~\% et 14.3~\%).
\textbf{12.} Corrélation réalisée de 0.03 à une seconde, 0.21 à dix, 0.49 à soixante et 0.59 à 300 secondes, pour une vraie corrélation de 0.6.
\textbf{13.} Hayashi--Yoshida 0.594 (écart type de 0.030 d'un jour à l'autre)~; \omterm{def:qm:high-frequency-econometrics:hy}{échantillonnage en temps de rafraîchissement} 0.48.
\textbf{14.} À cinq minutes, la RV atteint 1.43 fois la \omterm{def:qm:volatility-models:rv}{variance intégrée} et la \omterm{def:qm:high-frequency-econometrics:bipower}{variation bipuissance} 1.16 (contre 1.00 et 1.00), pour un saut qui représente 29~\% de la \omterm{def:qm:brownian-motion:qv}{variation quadratique}
de la journée.
\textbf{15.} La \omterm{def:qm:high-frequency-econometrics:bipower}{variation bipuissance} exige un échantillonnage fin pour isoler le saut, et un échantillonnage fin fait revenir le bruit~; les \omterm{def:qm:estimation:estimator}{estimateurs} robustes au bruit, en retour, comptent les sauts comme de la variance.
\textbf{16.} Une estimation robuste au bruit (\omterm{def:qm:high-frequency-econometrics:kernel}{noyau réalisé} ou pré-moyennage) sur la dernière fenêtre, combinée à une prévision quotidienne (chapitre~18).
\textbf{17.} Hayashi--Yoshida ou un \omterm{def:qm:estimation:estimator}{estimateur} synchronisé robuste au bruit, jamais des rendements à une seconde au dernier tick.
\textbf{18.} Un bruit qui dépend du signe de la transaction et de la fourchette, une négociation regroupée, et une volatilité variable au cours de la journée.
\textbf{19.} Résultat nommé~: \emph{la volatilité qui dépendait de l'horloge}~: avec un ratio bruit sur signal de $1.24 \times 10^{-4}$, l'intervalle d'échantillonnage optimal est d'environ 90 secondes (92
d'après le bruit estimé, 82 d'après le vrai bruit), et il évite une surestimation de la variance d'un facteur 5.8 par la \omterm{def:qm:volatility-models:rv}{variance réalisée} à une seconde (volatilité de 60~\% pour une vraie volatilité de 25~\%).
\textbf{20.} Parce que chaque observation supplémentaire ajoute la même variance de bruit alors que la variation efficiente qu'elle apporte diminue, si bien qu'au-delà d'une certaine fréquence la somme mesure le bruit.
\end{solution}

\begin{solution}{iq:qm:high-frequency-econometrics:1}
\omterm{def:qm:high-frequency-econometrics:noise}{Bruit de microstructure}~: le \omterm{def:qm:high-frequency-econometrics:noise}{rebond bid--ask} et la discrétisation des prix ajoutent une différence de bruit à chaque rendement, dont la variance ne diminue pas avec l'intervalle, si bien que la \omterm{def:qm:volatility-models:rv}{variance réalisée}
à partir de données tick compte $2n$ variances de bruit. L'échantillonnage à cinq minutes rend $n$ assez petit pour que le bruit soit négligeable.
\iqlookfor{Le biais $2n\omega^2$ et le \omterm{def:qm:high-frequency-econometrics:signature}{signature plot}.}
\end{solution}

\begin{solution}{iq:qm:high-frequency-econometrics:2}
\omterm{def:qm:high-frequency-econometrics:roll}{Estimateur de Roll}~: sous rebond, les variations consécutives des prix de transaction sont corrélées négativement, avec une covariance de $-c^2/4$, si bien que la fourchette effective vaut $2\sqrt{-\gamma_1}$. Avec les cotations et les
signes des transactions, comparer plutôt les prix de transaction au milieu en vigueur.
\iqlookfor{Le modèle de rebond et ses hypothèses.}
\end{solution}

\begin{solution}{iq:qm:high-frequency-econometrics:3}
L'intervalle qui équilibre la variance de discrétisation ($2\,\mathrm{IQ}/n$) et le biais de bruit ($2n\omega^2$)~: $n^* = (\mathrm{IQ}/4\omega^4)^{1/3}$. Il dépend du ratio bruit sur signal
ratio~: les journées plus volatiles et les grandes actions (grille relative plus fine) permettent un échantillonnage plus rapide.
\iqlookfor{Le compromis et la formule de la racine cubique.}
\end{solution}

\begin{solution}{iq:qm:high-frequency-econometrics:4}
L'\omterm{def:qm:high-frequency-econometrics:epps}{effet Epps}~: les deux actions ne négocient pas aux mêmes instants, si bien que sur une grille d'une seconde la plupart des rendements de l'une sont appariés à des rendements nuls de l'autre. Utiliser Hayashi--Yoshida,
un \omterm{def:qm:high-frequency-econometrics:hy}{échantillonnage en temps de rafraîchissement} avec un \omterm{def:qm:estimation:estimator}{estimateur} robuste au bruit, ou un échantillonnage plus grossier.
\iqlookfor{L'asynchronicité comme cause, et un \omterm{def:qm:estimation:estimator}{estimateur} de covariance qui n'interpole pas.}
\end{solution}

\begin{solution}{iq:qm:high-frequency-econometrics:5}
Comparer la \omterm{def:qm:volatility-models:rv}{variance réalisée}, qui inclut les sauts au carré, à la \omterm{def:qm:high-frequency-econometrics:bipower}{variation bipuissance} (ou à seuil), qui n'estime que la partie diffusive~; leur différence, standardisée,
est un test de saut. Utiliser un échantillonnage modéré ou des versions robustes au bruit, car le bruit gonfle les deux.
\iqlookfor{RV contre bipuissance et la mise en garde sur le bruit.}
\end{solution}

\begin{solution}{iq:qm:high-frequency-econometrics:6}
Un bruit iid rend les rendements observés MA(1)~: leur variance est gonflée de $2\omega^2$ et leur première autocovariance vaut $-\omega^2$. Ajouter deux fois les autocovariances pondérées à la variance
annule ce gonflement, exactement comme une variance de long terme somme les autocovariances d'un MA~; des pondérations lisses comme celle de Parzen maintiennent l'estimation positive et traitent le bruit dépendant.
\iqlookfor{La structure MA(1) et l'analogie HAC.}
\end{solution}
